\chapter{Grandfather Fights an Ostrich}
Before grandfather joined the Indian Railways, he worked for some time on the East African Railways, and it was during that period that he had his famous encounter with the ostrich. My childhood was frequently enlivened by this oft-told tale of my grandfather's, and I give it here in his own words-or as well as I can remember them:
\vspace{0.1in}
\par
While engaged in the laying of a new railway line, I had a miraculous escape from an awful death. I live in a small township, but my work lay some twelve miles away, and although I had a tent on the works, I often had to go into town on horseback.
On one occasion, an accident happening to my horse, I got a lift into town, hoping that someone might do me a similar favour on my way back. But this was not to be, and I made up my mind next morning to make the journey on foot, shortening the distance by talking a cut through the hills, which would save me about six miles.
\par
To take this shortcut, it was necessary to cross an ostrich 'camp' or farm. To venture across these 'camps' in the breeding season, especially on foot, can be dangerous, for during this time the male birds are extremely ferocious.
\par
But being familiar with the ways of ostriches, I knew my dog would scare away any ostrich that tried to attack me. Strange through it may seem, even the biggest ostrich (and some of them grow to a height of nine feet) will bolt faster than a racehorse at the same of even a small dog. And so, in company with my dog (a mongrel who had adopted me the previous month), I felt reasonably safe.
\par
On arrival at the 'camp', I got through the wire fencing and, keeping a good lookout, dodged across the spaces between the thorn bushes, now and then getting a sight of the birds that were feeding some distance away.
\par
I had gone about half a mile from the fencing when up started a hare, and in an instant my dog gave chase. I tried to call him back although I knew it was useless, since chasing hares was a passion with him.
\par
Whether it was the dog's bark on my own shouting, I don't know, but just what I was more anxious to avoid immediately happened: the ostriches were startled and began darting to and fro. Suddenly, I saw a big male bird emerge from a thicket about a hundred yards away. He stood still and stared at me for a few moments; then, expanding his wings and with his tail erect, he came bounding towards me.
\par
Believing discretion to be the better part of valour (at least in that particular situation), I turned and ran towards the fence. But it was an unequal race. What were my steps of two or three feet against the creature's great strides of sixteen to twenty feet? There was only one hope: to wait for the ostrich behind some bush and try to dodge it till it tried. A dodging game was obviously my only chance.
\par
Altering course a little, I rushed for the nearest clump of bushes where, gasping for breath, I waited for my pursuer. The great bird was almost immediately upon me, and a strange encounter commenced. This way and that, I dodged, talking great care that I did not get directly in front of its deadly kick. The ostrich kicks forward, and with such terrific force that his great chisel-like nails, if they stuck, would rip one open from head to foot.
\par
Breathless, and really quite helpless, I prayed wildly for help as I circled the bush, which was about twelve feet in diameter and some six feet in height. My strength was rapidly failing, and I realized it would be impossible to keep up the struggle much longer; I was ready to drop from sheer exhaustion. As if aware of my condition, the infuriated bird at me. With a desperate effort, I managed to step to one side. How it happened I don't know, but I found myself holding on to one of the creature's wings, close to its body.
\par
It was now the bird's turn to be frightened, and it began to turn, or rather waltz, moving round and round so quickly that my feet were soon swinging out almost horizontally. All the time the ostrich kept opening and shutting its break with loud snaps.
\par
Imagine my situation as I clung desperately to the wing of the enraged bird, which is whirling me round and round as if I had been cork! My arms soon began to ache with the strain, and the swift and continuous circling was making me dizzy. But I knew that if I relaxed my hold, a terrible fate awaited me: I should be promptly trampled to death by the spiteful bird.
\par
Round and round we went in a great circle. It seemed as if my enemy would never tire. But I knew I could not hold on much longer.
\par
Suddenly, the bird went into reverse! This unexpected movement not only had the effect of making me lose my hold but sent me sprawling to the ground. I landed in a heap at the foot of the thorn bush. In an instant, almost before I had time to realize what had happened, the ostrich was upon me. I thought the end had come. Instinctively, I put up my hands to protect my face. But, to my amazement, the great bird did not strike.
\par
I moved my hands from my face, and there stood the ostrich with one foot raised, ready to rip me open! I couldn't move. Was the bird going to play with me like a cat with a mouse, and prolong the agony?
\par
As I watched fascinated, I saw it turn its head sharply to the left. A second later, it jumped back, turned and made off as fast as it could go. Dazed, I wondered what had happened.
\par
I soon found out, for, to my great joy, I heard the bark of my truant dog, and the next moment he was jumping around me, licking my face and hands.
\par
Needless to say, I returned his caresses most affectionately! And I took good care to see that he did not leave my side until we were well clear of the ostrich 'camp'.
